%% kapa: exam=1 round=19 year=2008 period=1 form=A subject=civil_law subject_ko=민법 seq=8 number=8
%% source: Q-Net 감정평가사 기출문제 (cst00309) artlSeq=1001962 file=extracted/1008/1001962/2008년 제19회 감정평가사 1차시험 기출문제/1교시.hwp
%% conversion: kordoc 4.24.51 -> tools/gen1.py (auto); review=proofread-sonnet; flags=none
\begin{question}{8}
배우자 乙, 자녀 丙, 부친 丁을 두고 있는 甲은 해외출장을 가던 중 2006년 5월 1일 비행기가 추락하여 실종되었으며 지금까지 아무런 소식이 없다. 甲에 대한 실종선고에 관한 설명으로 틀린 것은? (다툼이 있는 경우에는 판례에 의함)
\begin{choices}
  \item 丙은 법률상 이해관계인이므로 실종선고를 청구할 수 있다.
  \item 丁은 재산상속을 받기 위해 실종선고를 청구할 수 없다.
  \item 乙이 실종선고를 신청하여 2007년 10월 1일에 실종선고를 받았다면 甲은 다음 날 오전 0시부터 사망한 것으로 간주된다.
  \item 실종선고에 의하여 甲이 권리능력을 상실하는 것은 아니다.
  \item 실종선고의 청구를 받은 가정법원은 6개월 이상 공시최고를 하여야 한다.
\end{choices}
\end{question}
\answer{3}
